\documentclass[11pt,a4paper]{article}
\usepackage[utf8]{inputenc}
\usepackage[T1]{fontenc}
\usepackage[catalan]{babel}
\usepackage{lmodern}
\usepackage[sfdefault,lining]{FiraSans}
\usepackage[a4paper,top=2.3cm,bottom=2.2cm,left=1.8cm,right=1.8cm,headheight=1.3cm,headsep=0.4cm,footskip=1cm]{geometry}
\usepackage{amsmath,amssymb}
\usepackage{icomma}
\usepackage{graphicx}
\usepackage[table]{xcolor}
\usepackage{tikz}
\usepackage{fontawesome5}
\usepackage{ccicons}
\usepackage{fancyhdr}
\usepackage{lastpage}
\usepackage{tasks}
\usepackage{needspace}
\usepackage[most]{tcolorbox}
\usepackage{array}
\usepackage[hidelinks,pdftitle={Fitxa d'exercicis de la unitat 2: Nombres reals. 3r d'ESO},pdfsubject={Matemàtiques 3r ESO},pdfkeywords={nombres reals, potències, radicals, notació científica, intervals}]{hyperref}

\definecolor{navy}{HTML}{1B2748}
\definecolor{accent}{HTML}{E8603A}
\definecolor{teal}{HTML}{1F9E90}
\definecolor{soft}{HTML}{EEF1F8}
\definecolor{muted}{HTML}{5D6685}

\setlength{\parindent}{0pt}
\setlength{\parskip}{3pt}
\settasks{label=\alph*), label-format=\bfseries\color{navy}, label-width=1.3em, item-indent=1.8em, column-sep=0.6em, after-item-skip=2pt, before-skip=3pt, after-skip=4pt}

% ---------- capçalera i peu ----------
\newcommand{\llicencia}{\href{https://creativecommons.org/licenses/by-nc-sa/4.0/deed.ca}{\ccbyncsa}\ \ Llicència Creative Commons Reconeixement-NoComercial-CompartirIgual 4.0 Internacional (CC BY-NC-SA 4.0)}
\pagestyle{fancy}
\fancyhf{}
\fancyhead[L]{\raisebox{-0.35\height}{\includegraphics[height=0.95cm]{logo-gva.png}}}
\fancyhead[C]{\small\color{muted}Fitxa d'exercicis · Unitat 2: Nombres reals · 3r d'ESO}
\fancyhead[R]{\raisebox{-0.35\height}{\includegraphics[height=0.95cm]{logo-vermellar.png}}}
\fancyfoot[L]{\scriptsize\color{muted}\llicencia}
\fancyfoot[R]{\small\color{muted}\thepage}
\renewcommand{\headrulewidth}{0.4pt}
\renewcommand{\headrule}{\color{muted!40}\hrule height \headrulewidth}
\fancypagestyle{primera}{\fancyhf{}\renewcommand{\headrulewidth}{0pt}
  \fancyfoot[L]{\scriptsize\color{muted}\llicencia}
  \fancyfoot[R]{\small\color{muted}\thepage}}

% ---------- nivells i exercicis ----------
\newcommand{\estrelles}[1]{{\color{accent}\ifcase#1\or\faStar\faStar[regular]\faStar[regular]\or\faStar\faStar\faStar[regular]\or\faStar\faStar\faStar\fi}}
\newcounter{exer}
\newcommand{\ex}[1]{\par\needspace{4\baselineskip}\medskip\stepcounter{exer}\noindent{\color{navy}\bfseries\theexer.}\ \ {\small\estrelles{#1}}\ \ }

\newcommand{\seccio}[2]{\par\needspace{10\baselineskip}\bigskip
  \begin{tcolorbox}[colback=navy,colframe=navy,coltext=white,arc=2mm,boxrule=0pt,left=3mm,right=3mm,top=1.2mm,bottom=1.2mm]
  {\large\bfseries #1\quad #2}
  \end{tcolorbox}\vspace{-2pt}}

% ---------- recta real ----------
\newcommand{\rectabuida}{\begin{tikzpicture}[x=0.42cm,baseline=0]
  \draw[-latex,muted] (-6.7,0)--(6.8,0);
  \foreach \i in {-6,...,6}{\draw[muted] (\i,0.12)--(\i,-0.12) node[below=-1pt,font=\scriptsize,text=muted]{$\i$};}
  \end{tikzpicture}}
% \rectaint{x1}{x2}{estil esquerra}{estil dreta}: estil = ple, buit, inf
\newcommand{\extrem}[2]{\ifx#2p\fill[accent] (#1,0) circle (0.13);\else\ifx#2b\filldraw[fill=white,draw=accent,line width=1.2pt] (#1,0) circle (0.13);\fi\fi}
\newcommand{\rectaint}[4]{\begin{tikzpicture}[x=0.42cm,baseline=0]
  \draw[-latex,muted] (-6.7,0)--(6.8,0);
  \foreach \i in {-6,...,6}{\draw[muted] (\i,0.12)--(\i,-0.12) node[below=-1pt,font=\scriptsize,text=muted]{$\i$};}
  \draw[accent,line width=3pt] (#1,0)--(#2,0);
  \ifx#3i\fill[accent] (#1-0.25,0)--(#1+0.15,0.2)--(#1+0.15,-0.2)--cycle;\else\extrem{#1}{#3}\fi
  \ifx#4i\fill[accent] (#2+0.25,0)--(#2-0.15,0.2)--(#2-0.15,-0.2)--cycle;\else\extrem{#2}{#4}\fi
  \end{tikzpicture}}

\newcolumntype{C}[1]{>{\centering\arraybackslash}p{#1}}

% =====================================================================
% Capçalera FentMates de les fitxes: un color per a cada curs.
% Els colors són els de la pàgina inicial (LEVELS a assets/js/materials.js);
% el logo ve de _eines/logo/pdf/ (python3 _eines/logo/genera.py).
%
% Ús, al final del preàmbul de la unitat:
%   % =====================================================================
% Capçalera FentMates de les fitxes: un color per a cada curs.
% Els colors són els de la pàgina inicial (LEVELS a assets/js/materials.js);
% el logo ve de _eines/logo/pdf/ (python3 _eines/logo/genera.py).
%
% Ús, al final del preàmbul de la unitat:
%   % =====================================================================
% Capçalera FentMates de les fitxes: un color per a cada curs.
% Els colors són els de la pàgina inicial (LEVELS a assets/js/materials.js);
% el logo ve de _eines/logo/pdf/ (python3 _eines/logo/genera.py).
%
% Ús, al final del preàmbul de la unitat:
%   \input{../comu/fentmates.tex}
%   \fmcurs{1eso}
% i, al començament del document:
%   \capcalera{Matemàtiques}{1r ESO}{Fitxa d'exercicis · Unitat 2}{Divisibilitat}
% Opcional abans de \input: \def\fmneed{14} (espai mínim abans de cada secció).
% =====================================================================
\makeatletter
\def\fm@logodir{../../logo/pdf}
\providecommand{\fmneed}{10}
\def\fm@cursnom{}

% clau del curs, color i color del text damunt del color
\newcommand{\fm@curs}[3]{\@namedef{fm@color@#1}{#2}\@namedef{fm@text@#1}{#3}}
\fm@curs{1eso}{E6B84A}{navy}
\fm@curs{2eso}{3B82F6}{white}
\fm@curs{3eso}{FF7A50}{white}
\fm@curs{4eso-a}{A78BFA}{white}
\fm@curs{4eso-b}{6D28D9}{white}
\fm@curs{1batx-mi}{9B1C1C}{white}
\fm@curs{1batx-ccss}{EF6B6B}{white}
\fm@curs{2batx-mii}{1E6B47}{white}
\fm@curs{2batx-ccss}{2BB3A3}{white}

\newcommand{\fmcurs}[1]{%
  \@ifundefined{fm@color@#1}{\PackageError{fentmates}{Curs desconegut: #1}{}}{}%
  \edef\fm@tmp{\noexpand\definecolor{curs}{HTML}{\@nameuse{fm@color@#1}}%
               \noexpand\colorlet{curstext}{\@nameuse{fm@text@#1}}}\fm@tmp
  \colorlet{cursfosc}{curs!72!black}%
  \def\fm@logo{\fm@logodir/vertical-#1.pdf}%
  \def\fm@marca{\fm@logodir/marca-#1.pdf}}

% ---------- capçalera de la primera pàgina ----------
% #1 matèria (petit), #2 curs, #3 tipus de document i unitat, #4 títol
\newcommand{\capcalera}[4]{%
  \gdef\fm@cursnom{#2}%
  \par\noindent
  \begin{tcolorbox}[colback=curs,colframe=curs,boxrule=0pt,arc=2.2mm,width=3.5cm,height=2.15cm,
      valign=center,left=3mm,right=2mm,top=1mm,bottom=1mm,nobeforeafter,box align=center]
    \color{curstext}{\footnotesize #1}\par\vspace{1pt}{\fontsize{19}{21}\selectfont\bfseries #2\par}%
  \end{tcolorbox}\hspace{5mm}%
  \begin{minipage}[c]{\dimexpr\linewidth-3.5cm-5mm-2.9cm\relax}
    {\color{cursfosc}\small\bfseries\MakeUppercase{#3}}\par\vspace{2pt}
    {\color{navy}\fontsize{20}{23}\selectfont\bfseries #4\par}
  \end{minipage}\hfill
  \begin{minipage}[c]{2.6cm}\raggedleft\includegraphics[width=2.6cm]{\fm@logo}\end{minipage}\par
  \vspace{2pt}}

% ---------- barres de secció del color del curs ----------
\renewcommand{\seccio}[2]{\par\needspace{\fmneed\baselineskip}\bigskip
  \begin{tcolorbox}[colback=curs,colframe=curs,coltext=curstext,arc=2mm,boxrule=0pt,left=3mm,right=3mm,top=1.2mm,bottom=1.2mm]
  {\large\bfseries #1\quad #2}
  \end{tcolorbox}\vspace{-2pt}}

% ---------- capçalera i peu de les pàgines ----------
% primera pàgina: logos institucionals al peu
\fancypagestyle{primera}{\fancyhf{}\renewcommand{\headrulewidth}{0pt}
  \fancyfoot[L]{\raisebox{-0.32\height}{\includegraphics[height=0.78cm]{logo-gva.png}}\hspace{5mm}%
    \raisebox{-0.32\height}{\includegraphics[height=0.78cm]{logo-vermellar.png}}}
  \fancyfoot[R]{\scriptsize\color{muted}\href{https://creativecommons.org/licenses/by-nc-sa/4.0/deed.ca}{\ccbyncsa}\ \ CC BY-NC-SA 4.0 · fentmates.org\qquad\small\thepage}}
% la resta: el símbol a l'esquerra i el curs a la dreta, amb el títol al centre com abans
\fancyhead[L]{\raisebox{-0.3\height}{\includegraphics[height=0.62cm]{\fm@marca}}}
\fancyhead[R]{\small\bfseries\color{cursfosc}\fm@cursnom}
\renewcommand{\headrule}{\color{curs!45}\hrule height \headrulewidth}
\makeatother

%   \fmcurs{1eso}
% i, al començament del document:
%   \capcalera{Matemàtiques}{1r ESO}{Fitxa d'exercicis · Unitat 2}{Divisibilitat}
% Opcional abans de \input: \def\fmneed{14} (espai mínim abans de cada secció).
% =====================================================================
\makeatletter
\def\fm@logodir{../../logo/pdf}
\providecommand{\fmneed}{10}
\def\fm@cursnom{}

% clau del curs, color i color del text damunt del color
\newcommand{\fm@curs}[3]{\@namedef{fm@color@#1}{#2}\@namedef{fm@text@#1}{#3}}
\fm@curs{1eso}{E6B84A}{navy}
\fm@curs{2eso}{3B82F6}{white}
\fm@curs{3eso}{FF7A50}{white}
\fm@curs{4eso-a}{A78BFA}{white}
\fm@curs{4eso-b}{6D28D9}{white}
\fm@curs{1batx-mi}{9B1C1C}{white}
\fm@curs{1batx-ccss}{EF6B6B}{white}
\fm@curs{2batx-mii}{1E6B47}{white}
\fm@curs{2batx-ccss}{2BB3A3}{white}

\newcommand{\fmcurs}[1]{%
  \@ifundefined{fm@color@#1}{\PackageError{fentmates}{Curs desconegut: #1}{}}{}%
  \edef\fm@tmp{\noexpand\definecolor{curs}{HTML}{\@nameuse{fm@color@#1}}%
               \noexpand\colorlet{curstext}{\@nameuse{fm@text@#1}}}\fm@tmp
  \colorlet{cursfosc}{curs!72!black}%
  \def\fm@logo{\fm@logodir/vertical-#1.pdf}%
  \def\fm@marca{\fm@logodir/marca-#1.pdf}}

% ---------- capçalera de la primera pàgina ----------
% #1 matèria (petit), #2 curs, #3 tipus de document i unitat, #4 títol
\newcommand{\capcalera}[4]{%
  \gdef\fm@cursnom{#2}%
  \par\noindent
  \begin{tcolorbox}[colback=curs,colframe=curs,boxrule=0pt,arc=2.2mm,width=3.5cm,height=2.15cm,
      valign=center,left=3mm,right=2mm,top=1mm,bottom=1mm,nobeforeafter,box align=center]
    \color{curstext}{\footnotesize #1}\par\vspace{1pt}{\fontsize{19}{21}\selectfont\bfseries #2\par}%
  \end{tcolorbox}\hspace{5mm}%
  \begin{minipage}[c]{\dimexpr\linewidth-3.5cm-5mm-2.9cm\relax}
    {\color{cursfosc}\small\bfseries\MakeUppercase{#3}}\par\vspace{2pt}
    {\color{navy}\fontsize{20}{23}\selectfont\bfseries #4\par}
  \end{minipage}\hfill
  \begin{minipage}[c]{2.6cm}\raggedleft\includegraphics[width=2.6cm]{\fm@logo}\end{minipage}\par
  \vspace{2pt}}

% ---------- barres de secció del color del curs ----------
\renewcommand{\seccio}[2]{\par\needspace{\fmneed\baselineskip}\bigskip
  \begin{tcolorbox}[colback=curs,colframe=curs,coltext=curstext,arc=2mm,boxrule=0pt,left=3mm,right=3mm,top=1.2mm,bottom=1.2mm]
  {\large\bfseries #1\quad #2}
  \end{tcolorbox}\vspace{-2pt}}

% ---------- capçalera i peu de les pàgines ----------
% primera pàgina: logos institucionals al peu
\fancypagestyle{primera}{\fancyhf{}\renewcommand{\headrulewidth}{0pt}
  \fancyfoot[L]{\raisebox{-0.32\height}{\includegraphics[height=0.78cm]{logo-gva.png}}\hspace{5mm}%
    \raisebox{-0.32\height}{\includegraphics[height=0.78cm]{logo-vermellar.png}}}
  \fancyfoot[R]{\scriptsize\color{muted}\href{https://creativecommons.org/licenses/by-nc-sa/4.0/deed.ca}{\ccbyncsa}\ \ CC BY-NC-SA 4.0 · fentmates.org\qquad\small\thepage}}
% la resta: el símbol a l'esquerra i el curs a la dreta, amb el títol al centre com abans
\fancyhead[L]{\raisebox{-0.3\height}{\includegraphics[height=0.62cm]{\fm@marca}}}
\fancyhead[R]{\small\bfseries\color{cursfosc}\fm@cursnom}
\renewcommand{\headrule}{\color{curs!45}\hrule height \headrulewidth}
\makeatother

%   \fmcurs{1eso}
% i, al començament del document:
%   \capcalera{Matemàtiques}{1r ESO}{Fitxa d'exercicis · Unitat 2}{Divisibilitat}
% Opcional abans de \input: \def\fmneed{14} (espai mínim abans de cada secció).
% =====================================================================
\makeatletter
\def\fm@logodir{../../logo/pdf}
\providecommand{\fmneed}{10}
\def\fm@cursnom{}

% clau del curs, color i color del text damunt del color
\newcommand{\fm@curs}[3]{\@namedef{fm@color@#1}{#2}\@namedef{fm@text@#1}{#3}}
\fm@curs{1eso}{E6B84A}{navy}
\fm@curs{2eso}{3B82F6}{white}
\fm@curs{3eso}{FF7A50}{white}
\fm@curs{4eso-a}{A78BFA}{white}
\fm@curs{4eso-b}{6D28D9}{white}
\fm@curs{1batx-mi}{9B1C1C}{white}
\fm@curs{1batx-ccss}{EF6B6B}{white}
\fm@curs{2batx-mii}{1E6B47}{white}
\fm@curs{2batx-ccss}{2BB3A3}{white}

\newcommand{\fmcurs}[1]{%
  \@ifundefined{fm@color@#1}{\PackageError{fentmates}{Curs desconegut: #1}{}}{}%
  \edef\fm@tmp{\noexpand\definecolor{curs}{HTML}{\@nameuse{fm@color@#1}}%
               \noexpand\colorlet{curstext}{\@nameuse{fm@text@#1}}}\fm@tmp
  \colorlet{cursfosc}{curs!72!black}%
  \def\fm@logo{\fm@logodir/vertical-#1.pdf}%
  \def\fm@marca{\fm@logodir/marca-#1.pdf}}

% ---------- capçalera de la primera pàgina ----------
% #1 matèria (petit), #2 curs, #3 tipus de document i unitat, #4 títol
\newcommand{\capcalera}[4]{%
  \gdef\fm@cursnom{#2}%
  \par\noindent
  \begin{tcolorbox}[colback=curs,colframe=curs,boxrule=0pt,arc=2.2mm,width=3.5cm,height=2.15cm,
      valign=center,left=3mm,right=2mm,top=1mm,bottom=1mm,nobeforeafter,box align=center]
    \color{curstext}{\footnotesize #1}\par\vspace{1pt}{\fontsize{19}{21}\selectfont\bfseries #2\par}%
  \end{tcolorbox}\hspace{5mm}%
  \begin{minipage}[c]{\dimexpr\linewidth-3.5cm-5mm-2.9cm\relax}
    {\color{cursfosc}\small\bfseries\MakeUppercase{#3}}\par\vspace{2pt}
    {\color{navy}\fontsize{20}{23}\selectfont\bfseries #4\par}
  \end{minipage}\hfill
  \begin{minipage}[c]{2.6cm}\raggedleft\includegraphics[width=2.6cm]{\fm@logo}\end{minipage}\par
  \vspace{2pt}}

% ---------- barres de secció del color del curs ----------
\renewcommand{\seccio}[2]{\par\needspace{\fmneed\baselineskip}\bigskip
  \begin{tcolorbox}[colback=curs,colframe=curs,coltext=curstext,arc=2mm,boxrule=0pt,left=3mm,right=3mm,top=1.2mm,bottom=1.2mm]
  {\large\bfseries #1\quad #2}
  \end{tcolorbox}\vspace{-2pt}}

% ---------- capçalera i peu de les pàgines ----------
% primera pàgina: logos institucionals al peu
\fancypagestyle{primera}{\fancyhf{}\renewcommand{\headrulewidth}{0pt}
  \fancyfoot[L]{\raisebox{-0.32\height}{\includegraphics[height=0.78cm]{logo-gva.png}}\hspace{5mm}%
    \raisebox{-0.32\height}{\includegraphics[height=0.78cm]{logo-vermellar.png}}}
  \fancyfoot[R]{\scriptsize\color{muted}\href{https://creativecommons.org/licenses/by-nc-sa/4.0/deed.ca}{\ccbyncsa}\ \ CC BY-NC-SA 4.0 · fentmates.org\qquad\small\thepage}}
% la resta: el símbol a l'esquerra i el curs a la dreta, amb el títol al centre com abans
\fancyhead[L]{\raisebox{-0.3\height}{\includegraphics[height=0.62cm]{\fm@marca}}}
\fancyhead[R]{\small\bfseries\color{cursfosc}\fm@cursnom}
\renewcommand{\headrule}{\color{curs!45}\hrule height \headrulewidth}
\makeatother

\fmcurs{3eso}

\begin{document}
\thispagestyle{primera}

% ---------- encapçalament ----------
\capcalera{Matemàtiques}{3r ESO}{Fitxa d'exercicis · Unitat 2}{Nombres reals}
\vspace{10pt}
Nom i cognoms: \rule{8.4cm}{0.4pt}\hfill Grup: \rule{1.6cm}{0.4pt}\hfill Data: \rule{2.2cm}{0.4pt}
\vspace{8pt}

\begin{tcolorbox}[colback=soft,colframe=soft,arc=2mm,left=3mm,right=3mm,top=1.5mm,bottom=1.5mm]
\small Exercicis graduats per nivells de dificultat:\quad
\estrelles{1}\ bàsic \qquad \estrelles{2}\ mitjà \qquad \estrelles{3}\ avançat.\\[2pt]
Resol-los al quadern mostrant tots els passos. Les rectes i les taules es poden completar directament en la fitxa.
\end{tcolorbox}

% =====================================================
\seccio{1}{Potències d'exponent natural}

\ex{1} Escriu com a potència i calcula:
\begin{tasks}(2)
  \task $5\cdot5\cdot5$
  \task $(-2)(-2)(-2)(-2)(-2)$
  \task $\frac23\cdot\frac23\cdot\frac23\cdot\frac23$
  \task $(-1)(-1)(-1)(-1)(-1)(-1)(-1)$
\end{tasks}

\ex{1} Calcula:
\begin{tasks}(4)
  \task $2^{6}$
  \task $(-3)^{3}$
  \task $(-3)^{4}$
  \task $-3^{4}$
  \task $(-1)^{100}$
  \task $\left(-\frac12\right)^{3}$
  \task $0^{9}$
  \task $10^{7}$
\end{tasks}

\ex{2} Sense fer cap càlcul, indica si el resultat és positiu o negatiu. Justifica-ho.
\begin{tasks}(4)
  \task $(-7)^{12}$
  \task $(-5)^{9}$
  \task $-(-2)^{6}$
  \task $-(-4)^{3}$
\end{tasks}

\ex{2} Calcula:
\begin{tasks}(2)
  \task $-2^{2}+(-2)^{2}-(-2)^{3}$
  \task $(-1)^{5}+(-1)^{6}-1^{7}$
  \task $3^{2}-2^{3}+(-1)^{10}$
  \task $\left(-\frac23\right)^{2}\cdot\left(\frac32\right)^{3}$
\end{tasks}

\ex{3} Troba el valor de $x$ en cada cas:
\begin{tasks}(4)
  \task $2^{x}=64$
  \task $(-3)^{x}=-27$
  \task $x^{4}=81$
  \task $\left(\frac25\right)^{x}=\frac{8}{125}$
\end{tasks}

% =====================================================
\seccio{2}{Propietats de les potències}

\ex{1} Expressa com una sola potència:
\begin{tasks}(3)
  \task $3^{4}\cdot3^{5}$
  \task $7^{10}:7^{6}$
  \task $\left(2^{3}\right)^{5}$
  \task $5^{3}\cdot5\cdot5^{4}$
  \task $a^{8}:a^{3}\ \ (a\neq0)$
  \task $\left((-2)^{2}\right)^{3}$
\end{tasks}

\ex{1} Aplica les propietats de les potències amb el mateix exponent i calcula:
\begin{tasks}(4)
  \task $2^{5}\cdot5^{5}$
  \task $12^{4}:4^{4}$
  \task $(-3)^{3}\cdot(-2)^{3}$
  \task $\left(\frac34\right)^{2}\cdot\left(\frac43\right)^{2}$
\end{tasks}

\ex{2} Simplifica i calcula:
\begin{tasks}(4)
  \task $\dfrac{2^{7}\cdot2^{3}}{\left(2^{4}\right)^{2}}$
  \task $\dfrac{\left(3^{2}\right)^{4}\cdot3^{3}}{3^{5}\cdot3^{4}}$
  \task $\dfrac{5^{6}\cdot5^{3}}{\left(5^{2}\right)^{4}}$
  \task $\dfrac{6^{3}\cdot6^{5}}{\left(6^{2}\right)^{3}}$
\end{tasks}

\ex{2} Descompon les bases en factors primers, simplifica i calcula:
\begin{tasks}(4)
  \task $\dfrac{4^{3}\cdot8^{2}}{16^{2}}$
  \task $\dfrac{9^{2}\cdot27}{3^{6}}$
  \task $\dfrac{25^{3}}{5^{2}\cdot125}$
  \task $\dfrac{6^{4}}{2^{3}\cdot3^{2}}$
\end{tasks}

\ex{3} Descompon, simplifica i calcula:
\begin{tasks}(4)
  \task $\dfrac{12^{3}\cdot18^{2}}{8^{2}\cdot27^{2}}$
  \task $\dfrac{15^{2}\cdot10^{3}}{6^{2}\cdot25^{2}}$
  \task $\dfrac{20^{3}\cdot9^{2}}{30^{2}\cdot4^{2}\cdot3}$
  \task $\dfrac{14^{2}\cdot21^{3}}{49^{2}\cdot6^{2}}$
\end{tasks}

\ex{3} Indica si cada igualtat és vertadera o falsa. Justifica la resposta.
\begin{tasks}(3)
  \task $2^{3}+2^{3}=2^{4}$
  \task $3^{2}\cdot3^{3}=9^{5}$
  \task $(2+3)^{2}=2^{2}+3^{2}$
  \task $5^{3}\cdot2^{3}=10^{3}$
  \task $2^{10}+2^{10}=2^{11}$
  \task $4^{5}:4^{5}=0$
\end{tasks}

% =====================================================
\seccio{3}{Potències d'exponent enter negatiu}

\ex{1} Calcula:
\begin{tasks}(4)
  \task $3^{-2}$
  \task $(-2)^{-3}$
  \task $10^{-3}$
  \task $\left(\frac14\right)^{-2}$
  \task $\left(\frac23\right)^{-3}$
  \task $\left(-\frac35\right)^{-2}$
  \task $7^{0}$
  \task $-2^{-4}$
\end{tasks}

\ex{1} Escriu com una potència d'exponent negatiu:
\begin{tasks}(5)
  \task $\dfrac{1}{5^{4}}$
  \task $\dfrac{1}{27}$
  \task $\dfrac{1}{16}$
  \task $0,001$
  \task $\dfrac{1}{49}$
\end{tasks}

\ex{2} Expressa com una sola potència:
\begin{tasks}(3)
  \task $5^{-3}\cdot5^{7}$
  \task $2^{-4}:2^{3}$
  \task $\left(3^{-2}\right)^{-3}$
  \task $7^{-2}\cdot7^{-5}:7^{-9}$
  \task $\left(4^{-1}\right)^{3}\cdot4^{5}$
  \task $10^{-3}\cdot10^{5}\cdot10^{-6}$
\end{tasks}

\ex{2} Calcula i dona el resultat en forma de fracció:
\begin{tasks}(4)
  \task $2^{-1}+2^{-2}$
  \task $3^{-1}-3^{-2}$
  \task $\left(\frac12\right)^{-2}+\left(\frac13\right)^{-1}$
  \task $5^{0}-5^{-1}$
\end{tasks}

\ex{3} Descompon, simplifica i calcula:
\begin{tasks}(3)
  \task $\dfrac{8^{-2}\cdot4^{3}}{2^{-5}\cdot16}$
  \task $\dfrac{9^{-1}\cdot27^{2}}{3^{-2}\cdot81}$
  \task $\dfrac{6^{-2}\cdot12^{3}}{18^{-1}\cdot2^{5}}$
  \task $\dfrac{10^{-2}\cdot25^{3}}{5^{4}\cdot2^{-3}}$
  \task $\left(\dfrac23\right)^{-2}\cdot\left(\dfrac94\right)^{-1}$
  \task $\left[\left(\dfrac12\right)^{-3}\cdot4^{-2}\right]^{-1}$
\end{tasks}

% =====================================================
\seccio{4}{Notació científica}

\ex{1} Escriu en notació científica:
\begin{tasks}(3)
  \task $45\,000\,000$
  \task $0,000\,72$
  \task $908\,000$
  \task $0,000\,000\,3$
  \task $-12\,500$
  \task $0,05$
\end{tasks}

\ex{1} Escriu en forma decimal:
\begin{tasks}(5)
  \task $3,2\cdot10^{6}$
  \task $5\cdot10^{-3}$
  \task $1,07\cdot10^{9}$
  \task $8,45\cdot10^{-6}$
  \task $-6,1\cdot10^{4}$
\end{tasks}

\ex{2} Indica quines d'aquestes expressions no estan escrites en notació científica i corregeix-les:
\begin{tasks}(3)
  \task $2,7\cdot10^{5}$
  \task $34\cdot10^{6}$
  \task $0,8\cdot10^{-3}$
  \task $10,5\cdot10^{-2}$
  \task $125\cdot10^{-8}$
  \task $0,0042\cdot10^{9}$
\end{tasks}

\ex{2} Ordena aquests nombres de menor a major:
\[
3,1\cdot10^{-4}\qquad 2,9\cdot10^{-3}\qquad 5\cdot10^{-5}\qquad 1,2\cdot10^{-4}\qquad 9\cdot10^{-6}
\]

\ex{3} Expressa en notació científica les dades següents:
\begin{tasks}(1)
  \task La velocitat de la llum és de $300\,000$ km/s.
  \task El diàmetre d'un glòbul roig és de $0,000\,007\,5$ m.
  \task La població mundial és d'uns 8\,000 milions de persones.
  \task La massa d'un àtom d'hidrogen és de $0,000\,000\,000\,000\,000\,000\,000\,001\,67$ g.
  \task Un any llum equival a uns $9\,460\,000\,000\,000$ km.
\end{tasks}

% =====================================================
\seccio{5}{Radicals. Extracció i introducció de factors}

\ex{1} Calcula, si existeixen en $\mathbb{R}$:
\begin{tasks}(4)
  \task $\sqrt{81}$
  \task $\sqrt[3]{-125}$
  \task $\sqrt[4]{16}$
  \task $\sqrt{-9}$
  \task $\sqrt[5]{-32}$
  \task $\sqrt{\dfrac{25}{49}}$
  \task $\sqrt[3]{\dfrac{8}{27}}$
  \task $\sqrt[6]{1}$
\end{tasks}

\ex{1} Extreu tots els factors possibles del radical:
\begin{tasks}(6)
  \task $\sqrt{12}$
  \task $\sqrt{18}$
  \task $\sqrt{45}$
  \task $\sqrt{98}$
  \task $\sqrt[3]{16}$
  \task $\sqrt[3]{81}$
\end{tasks}

\ex{1} Introdueix els factors dins del radical:
\begin{tasks}(6)
  \task $2\sqrt{5}$
  \task $3\sqrt{3}$
  \task $5\sqrt{2}$
  \task $2\sqrt[3]{3}$
  \task $3\sqrt[4]{2}$
  \task $a\sqrt{b}$
\end{tasks}

\ex{2} Extreu tots els factors possibles del radical:
\begin{tasks}(3)
  \task $\sqrt{200}$
  \task $\sqrt{2^{5}\cdot3^{3}}$
  \task $\sqrt[3]{2^{7}\cdot5^{4}}$
  \task $\sqrt[4]{3^{9}}$
  \task $\sqrt{a^{7}\,b^{4}}$
  \task $\sqrt[3]{x^{5}\,y^{9}}$
\end{tasks}

\ex{2} Introdueix els factors dins del radical i ordena de menor a major:
\[
3\sqrt{5}\qquad 4\sqrt{3}\qquad 2\sqrt{11}\qquad 5\sqrt{2}
\]

\ex{3} Extreu tots els factors possibles:
\begin{tasks}(4)
  \task $\sqrt{\dfrac{50}{81}}$
  \task $\sqrt{12a^{3}}$
  \task $\sqrt[3]{-54}$
  \task $\sqrt[4]{32\,a^{6}\,b^{5}}$
\end{tasks}

% =====================================================
\seccio{6}{Operacions amb radicals}

\ex{1} Opera aquests radicals semblants:
\begin{tasks}(2)
  \task $3\sqrt{5}+2\sqrt{5}$
  \task $7\sqrt{3}-4\sqrt{3}+\sqrt{3}$
  \task $2\sqrt[3]{6}-5\sqrt[3]{6}+8\sqrt[3]{6}$
  \task $\sqrt{2}+3\sqrt{2}-6\sqrt{2}$
\end{tasks}

\ex{1} Calcula i simplifica:
\begin{tasks}(6)
  \task $\sqrt{3}\cdot\sqrt{12}$
  \task $\sqrt{5}\cdot\sqrt{20}$
  \task $\sqrt{2}\cdot\sqrt{6}$
  \task $\sqrt[3]{2}\cdot\sqrt[3]{4}$
  \task $\sqrt{48}:\sqrt{3}$
  \task $\sqrt[3]{54}:\sqrt[3]{2}$
\end{tasks}

\ex{2} Extreu factors i opera:
\begin{tasks}(2)
  \task $\sqrt{8}+\sqrt{18}$
  \task $\sqrt{12}+\sqrt{75}-\sqrt{27}$
  \task $2\sqrt{20}-\sqrt{45}+\sqrt{80}$
  \task $\sqrt{50}-2\sqrt{32}+3\sqrt{2}$
  \task $\sqrt[3]{16}+\sqrt[3]{54}-\sqrt[3]{2}$
  \task $\sqrt{28}+\sqrt{63}-\sqrt{7}$
\end{tasks}

\ex{2} Calcula i simplifica:
\begin{tasks}(3)
  \task $2\sqrt{3}\cdot5\sqrt{6}$
  \task $\sqrt{10}\cdot\sqrt{15}$
  \task $\left(3\sqrt{2}\right)^{2}$
  \task $12\sqrt{30}:\left(4\sqrt{6}\right)$
  \task $\left(\sqrt{3}\right)^{5}$
  \task $\sqrt{2}\cdot\sqrt{3}\cdot\sqrt{6}$
\end{tasks}

\ex{3} Opera i simplifica:
\begin{tasks}(3)
  \task $\sqrt{2}\left(\sqrt{8}+\sqrt{18}\right)$
  \task $\sqrt{3}\left(2\sqrt{3}-\sqrt{12}\right)$
  \task $\left(\sqrt{6}+\sqrt{2}\right)\cdot\sqrt{2}$
  \task $\sqrt{12}+\sqrt{27}-\sqrt{3}\cdot\sqrt{6}$
  \task $\dfrac{\sqrt{50}-\sqrt{8}}{\sqrt{2}}$
\end{tasks}

% =====================================================
\seccio{7}{Intervals: definició i tipus}

\ex{1} Completa la taula:\par\smallskip
{\renewcommand{\arraystretch}{1.75}\centering
\begin{tabular}{|C{3.6cm}|C{4.2cm}|C{6.2cm}|}
\hline
\rowcolor{soft}\textbf{Interval} & \textbf{Desigualtat} & \textbf{Tipus d'interval} \\ \hline
$(-3,\ 5)$ & & \\ \hline
$[-1,\ 4]$ & & \\ \hline
 & $2<x\le 6$ & \\ \hline
 & $x\ge -2$ & \\ \hline
$(-\infty,\ 0)$ & & \\ \hline
 & $-5\le x<-1$ & \\ \hline
\end{tabular}\par}

\ex{1} Escriu com a interval i com a desigualtat:
\begin{tasks}(2)
  \task Els nombres majors que $-4$.
  \task Els nombres menors o iguals que $7$.
  \task Els nombres compresos entre $0$ i $10$, tots dos inclosos.
  \task Els nombres majors que $-2$ i menors o iguals que $3$.
\end{tasks}

\ex{2} Indica si el nombre pertany o no a l'interval:
\begin{tasks}(3)
  \task $3\in(-1,\ 3)$?
  \task $-2\in[-2,\ 5)$?
  \task $0\in(-\infty,\ 0]$?
  \task $2,5\in(2,\ 3)$?
  \task $-7\in(-6,\ +\infty)$?
  \task $\sqrt{2}\in[1,\ 2]$?
\end{tasks}

\ex{3} Respon:
\begin{tasks}(1)
  \task Escriu tres nombres que pertanyen a l'interval $(1,\ 2)$. Quants nombres reals conté aquest interval?
  \task Escriu un interval que continga el $0$ i el $5$, però no el $-1$ ni el $6$.
  \task Quin és el major nombre enter que pertany a $[-3,\ 4)$? I el major que pertany a $(-3,\ 4]$?
  \task Escriu com a interval el conjunt de nombres reals la distància dels quals al $2$ és menor que $3$.
\end{tasks}

% =====================================================
\seccio{8}{Representació gràfica d'intervals}

\ex{1} Representa a la recta real els intervals següents:\par\smallskip
\begin{tabular}{@{}l@{\ \ }l@{\qquad}l@{\ \ }l@{}}
a) $[-2,\ 3]$ & \rectabuida & b) $(0,\ 4)$ & \rectabuida \\[16pt]
c) $(-\infty,\ 1]$ & \rectabuida & d) $(-3,\ +\infty)$ & \rectabuida \\[16pt]
e) $[-4,\ -1)$ & \rectabuida & f) $x\le 2$ & \rectabuida
\end{tabular}

\ex{2} Escriu l'interval i la desigualtat que corresponen a cada representació:\par\smallskip
\begin{tabular}{@{}l@{\ \ }l@{\qquad}l@{\ \ }l@{}}
a) & \rectaint{-1}{4}{b}{p} & b) & \rectaint{2}{6.4}{p}{i} \\[16pt]
c) & \rectaint{-6.4}{-3}{i}{b} & d) & \rectaint{-4}{0}{p}{p}
\end{tabular}

\ex{3} Representa a la recta real i escriu com a interval (o unió d'intervals):\par\smallskip
\begin{tabular}{@{}l@{\ \ }l@{}}
a) $x>-1$\ \ i\ \ $x\le 5$ & \rectabuida \\[16pt]
b) $x\ge -3$\ \ i\ \ $x<2$ & \rectabuida \\[16pt]
c) $x\le -1$\ \ o\ \ $x>4$ & \rectabuida
\end{tabular}

% =====================================================
\seccio{9}{Unió i intersecció d'intervals}

\ex{1} Siguen $A=[-2,\ 4]$ i $B=(1,\ 6)$. Representa'ls en la mateixa recta i calcula $A\cup B$ i $A\cap B$.\par\smallskip
{\centering\rectabuida\par\vspace{12pt}\rectabuida\par}

\ex{2} Calcula $A\cup B$ i $A\cap B$ en cada cas (fes-ne abans la representació gràfica):
\begin{tasks}(2)
  \task $A=(-\infty,\ 3)$,\ \ $B=[0,\ +\infty)$
  \task $A=[-5,\ -1]$,\ \ $B=(-3,\ 2)$
  \task $A=(-\infty,\ -2]$,\ \ $B=(1,\ +\infty)$
  \task $A=[0,\ 5)$,\ \ $B=[5,\ 8]$
  \task $A=(-4,\ 4)$,\ \ $B=[-1,\ 1]$
\end{tasks}

\ex{3} Respon:
\begin{tasks}(1)
  \task Escriu dos intervals $A$ i $B$ tals que $A\cap B=[2,\ 3]$.
  \task Siguen $A=[-3,\ 5)$ i $B=(a,\ 8]$. Quin valor ha de tindre $a$ perquè $A\cap B=(1,\ 5)$?
  \task Siguen $A=[-2,\ 6]$, $B=(0,\ +\infty)$ i $C=(-\infty,\ 3]$. Calcula $A\cap B\cap C$ i $(A\cap B)\cup C$.
  \task Escriu l'interval $[-1,\ 7)$ com la intersecció de dues semirectes.
\end{tasks}

% =====================================================
\seccio{10}{Classificació dels nombres reals}

\ex{1} Marca amb una creu els conjunts als quals pertany cada nombre:\par\smallskip
{\renewcommand{\arraystretch}{1.3}\centering
\begin{tabular}{|C{3.6cm}|C{1.3cm}|C{1.3cm}|C{1.3cm}|C{1.3cm}|C{1.3cm}|}
\hline
\rowcolor{soft}\textbf{Nombre} & $\mathbb{N}$ & $\mathbb{Z}$ & $\mathbb{Q}$ & $\mathbb{I}$ & $\mathbb{R}$ \\ \hline
$12$ & & & & & \\ \hline
$-5$ & & & & & \\ \hline
$\dfrac37$ & & & & & \\ \hline
$0,25$ & & & & & \\ \hline
$\sqrt{3}$ & & & & & \\ \hline
$\pi$ & & & & & \\ \hline
$-\sqrt{49}$ & & & & & \\ \hline
$0,\overline{3}$ & & & & & \\ \hline
$1,234\,567\,891\,011\ldots$ & & & & & \\ \hline
\end{tabular}\par}

\ex{2} Indica el conjunt numèric més xicotet al qual pertany cada nombre:
\begin{tasks}(4)
  \task $\dfrac{18}{6}$
  \task $-\sqrt{25}$
  \task $\sqrt{8}$
  \task $2,\overline{5}$
  \task $1+\sqrt{2}$
  \task $-\dfrac{27}{9}$
  \task $\sqrt{0,09}$
  \task $0,101\,001\,000\,1\ldots$
  \task $\sqrt[3]{-64}$
  \task $\dfrac{\pi}{2}$
\end{tasks}

\ex{2} Indica si les afirmacions següents són vertaderes o falses. Justifica-ho.
\begin{tasks}(2)
  \task Tot nombre natural és enter.
  \task Tot nombre enter és natural.
  \task Tot nombre decimal periòdic és racional.
  \task L'arrel quadrada d'un natural sempre és irracional.
  \task Hi ha nombres que són racionals i irracionals alhora.
  \task Tots els nombres decimals són racionals.
\end{tasks}

\ex{3} Respon:
\begin{tasks}(1)
  \task Escriu un nombre irracional comprés entre $3$ i $4$.
  \task Escriu un nombre racional comprés entre $\sqrt{2}$ i $\sqrt{3}$.
  \task Per a quins valors naturals de $n$, entre $1$ i $30$, el nombre $\sqrt{n}$ és racional?
  \task La suma de dos nombres irracionals és sempre irracional? Busca'n un contraexemple.
\end{tasks}

\vfill
\begin{center}\color{muted}\small\faCheckCircle[regular]\ \ Ànim! Si et costa algun exercici, repassa els apunts interactius de la unitat.\end{center}

\end{document}
